	\documentclass[12pt,a4paper]{ctexart}
\usepackage{amsmath, amssymb, amsfonts}
\usepackage{geometry}
\usepackage{tcolorbox}
\usepackage{enumitem}
\usepackage{tikz}
\usepackage[european]{circuitikz}
\usepackage{pgfplots}
\pgfplotsset{compat=1.17}
% 修复：加载 positioning 库以支持 below=of 语法
\usetikzlibrary{positioning, arrows.meta, calc}
\geometry{left=2.2cm, right=2.2cm, top=2.5cm, bottom=2.5cm}
\tcbuselibrary{breakable, skins, theorems}
\definecolor{mainblue}{RGB}{0, 102, 204}
\definecolor{maingreen}{RGB}{0, 153, 76}
\definecolor{mainorange}{RGB}{255, 102, 0}
\definecolor{mainred}{RGB}{204, 0, 0}
\definecolor{lightgray}{RGB}{240, 240, 240}
\newtcolorbox{notebox}[1][]{
	colback=mainblue!5!white,
	colframe=mainblue!75!black,
	fonttitle=\bfseries,
	title=#1,
	breakable,
	enhanced jigsaw,
	arc=2mm,
	boxrule=1pt
}
\newtcolorbox{warnbox}[1][]{
	colback=mainorange!5!white,
	colframe=mainorange!75!black,
	fonttitle=\bfseries,
	title=#1,
	breakable,
	enhanced jigsaw,
	arc=2mm,
	boxrule=1pt
}
\newtcolorbox{tipbox}[1][]{
	colback=maingreen!5!white,
	colframe=maingreen!75!black,
	fonttitle=\bfseries,
	title=#1,
	breakable,
	enhanced jigsaw,
	arc=2mm,
	boxrule=1pt
}
\newtcolorbox{derivebox}[1][]{
	colback=lightgray,
	colframe=black!50,
	fonttitle=\bfseries,
	title=#1,
	breakable,
	enhanced jigsaw,
	arc=1mm,
	boxrule=0.5pt,
	left=5pt, right=5pt
}
% 修复：标题换行符使用 \\
\title{\textbf{《模拟电路与数字电路》\\ 第1$\sim$5章 极致复习笔记}}
\author{寇戈 \quad 蒋立平 \quad 吴少琴 \quad 编著\\ \small 智能问答专家精排版}
\date{当前基准时间：2026年5月14日}
\begin{document}
	\maketitle
	\thispagestyle{empty}
	\begin{center}
		\vspace{2cm}
		\Large\textbf{前 言}
	\end{center}
	模电被无数工科生戏称为“魔电”，其痛点在于微观物理机制抽象、电路模型繁多、交直流叠加分析复杂。本笔记基于寇戈、蒋立平、吴少琴编著的《模拟电路与数字电路》前五章，旨在通过\textbf{形象比喻}、\textbf{详尽推导}与\textbf{精美图解}，打破认知壁垒。
	\clearpage
	\tableofcontents
	\clearpage
	\section*{第1章 绪论}
	\addcontentsline{toc}{section}{第1章 绪论}
	\subsection*{1.1 电子信号与电子系统}
	\begin{notebox}[核心概念：模拟与数字的本质]
		\begin{itemize}[leftmargin=2em]
			\item \textbf{模拟信号}：时间上和幅度上均连续的信号。\\
			\textbf{形象比喻}：就像水流，连续不断，水位可以停留在任意高度。自然界中感知的信号（声、光、温）多为模拟信号。
			\item \textbf{数字信号}：时间上和幅度上均离散的信号，通常用0和1表征。\\
			\textbf{形象比喻}：就像楼梯，只能站在特定的台阶上，且跳动是瞬间完成的。计算机内部流动的全是数字信号。
		\end{itemize}
	\end{notebox}
	电子系统的核心任务往往是：感知模拟世界 $\to$ 转换为数字信号处理 $\to$ 转换回模拟信号执行。这其中，放大器是所有模拟系统的灵魂。
	\subsection*{1.2 放大的基本概念}
	放大电路的本质不是凭空产生能量，而是\textbf{用微弱的输入信号去控制电源的能量，使负载获得与输入信号成比例的大能量输出}。
	\begin{itemize}
		\item \textbf{电压增益}：$A_u = \frac{u_o}{u_i}$，常以分贝 表示：$A_u(\text{dB}) = 20\lg|A_u|$。
		\item \textbf{输入电阻}：$R_i = \frac{u_i}{i_i}$。相当于信号源的“负载”，越大越好，可减轻信号源负担（电压放大）。
		\item \textbf{输出电阻}：$R_o$。相当于放大器内部的“内阻”，越小越好，带载能力强（电压放大）。
	\end{itemize}
	\clearpage
	\section*{第2章 半导体器件基础}
	\addcontentsline{toc}{section}{第2章 半导体器件基础}
	本章是模电的基石，微观物理机制决定了宏观电路特性。
	\subsection*{2.1 半导体基础}
	\begin{itemize}
		\item \textbf{本征半导体}：纯净的半导体（如Si、Ge）。在绝对零度时为绝缘体；温度升高，本征激发产生\textbf{电子-空穴对}，载流子浓度急剧上升。
		\item \textbf{杂质半导体}：
		\begin{itemize}
			\item \textbf{N型}：掺入五价元素（如磷P），多子是\textbf{电子}，少子是空穴。
			\item \textbf{P型}：掺入三价元素（如硼B），多子是\textbf{空穴}，少子是电子。
		\end{itemize}
		\textbf{形象记忆}：N代表Negative（负），电子带负电；P代表Positive（正），空穴带正电。多子由掺杂决定，少子由温度决定（故少子对温度极敏感）。
	\end{itemize}
	\subsection*{2.2 PN结的形成与特性}
	PN结是半导体器件的“灵魂”，其形成是两种相反运动动态平衡的结果。
	\begin{tipbox}[形象理解PN结：两军对垒]
		将P区和N区想象成两支军队。P区全是“空穴兵”（正），N区全是“电子兵”（负）。相遇时，前线士兵互相同归于尽（扩散运动），留下无法移动的带电离子（空间电荷区/耗尽层），形成内电场。内电场如同战壕，阻止后方士兵继续冲锋（漂移运动），最终达到动态平衡。
	\end{tipbox}
	\begin{center}
		\begin{tikzpicture}[scale=0.9, >=stealth]
			% P region
			\fill[blue!10] (0,0) rectangle (3,4);
			\node at (1.5,4.3) {\textbf{P型区}};
			\foreach \x/\y in {0.5/0.5, 1.5/1.5, 2.5/0.8, 0.8/2.8, 2.2/3.2, 1/1, 2/2.5} {
				\draw[thick, red] (\x,\y) circle (0.15);
				\draw[thick, red] (\x-0.15,\y) -- (\x+0.15,\y);
			}
			% N region
			\fill[red!10] (3,0) rectangle (6,4);
			\node at (4.5,4.3) {\textbf{N型区}};
			\foreach \x/\y in {3.5/0.8, 4.5/1.8, 5.5/0.5, 3.8/2.5, 5.2/3.0, 4/1, 5/3.5} {
				\draw[thick, blue] (\x,\y) circle (0.15);
				\draw[thick, blue] (\x,\y-0.15) -- (\x,\y+0.15);
			}
			% Depletion region
			\fill[gray!30] (2.4,0) rectangle (3.6,4);
			\node at (3, -0.4) {\textbf{空间电荷区}};
			% Ions in depletion
			\foreach \y in {0.5, 1.5, 2.5, 3.5} {
				\draw[thick, red] (2.6,\y) circle (0.12);
				\draw[thick, red] (2.45,\y) -- (2.75,\y); % Negative ion in P boundary
			}
			\foreach \y in {0.8, 1.8, 2.8, 3.2} {
				\draw[thick, blue] (3.4,\y) circle (0.12);
				\draw[thick, blue] (3.4,\y-0.12) -- (3.4,\y+0.12); % Positive ion in N boundary
			}
			% E field
			\draw[->, ultra thick, orange] (3.4, 4.6) -- (2.6, 4.6) node[midway, above] {$E_{内}$};
		\end{tikzpicture}
	\end{center}
	\begin{warnbox}[PN结的单向导电性]
		\begin{itemize}
			\item \textbf{正偏（P接正，N接负）}：外电场削弱内电场，耗尽层变窄，多子扩散加剧，形成大电流，称为\textbf{导通}。犹如外力推平战壕，士兵疯狂冲锋。
			\item \textbf{反偏（P接负，N接正）}：外电场增强内电场，耗尽层变宽，少子漂移形成极小电流，称为\textbf{截止}。犹如加深战壕，只有极少数精锐（少子）能勉强通过。
		\end{itemize}
	\end{warnbox}
	\subsection*{2.3 半导体二极管}
	\subsubsection*{2.3.1 伏安特性与模型}
	二极管的本质就是一个PN结。其伏安特性方程为：
	$$ I = I_S(e^{U/U_T} - 1) $$
	其中，$I_S$ 为反向饱和电流，常温下热电压 $U_T \approx 26\text{mV}$。
	\begin{center}
		\begin{tikzpicture}
			\begin{axis}[
				width=12cm, height=7cm,
				xlabel={电压 $U$ / V},
				ylabel={电流 $I$ / mA},
				xtick={-5, -0.7, 0, 0.5, 0.7},
				ytick={0, 10, 20},
				xmin=-6, xmax=1,
				ymin=-1, ymax=25,
				axis lines=middle,
				thick,
				grid=major,
				grid style={dashed, gray!30},
				every axis x label/.style={at={(ticklabel* cs:1.0)}, anchor=west},
				every axis y label/.style={at={(ticklabel* cs:1.0)}, anchor=south},
				restrict y to domain=-5:30 % 修复：限制y轴范围防止指数溢出
				]
				% Forward characteristic (exponential)
				% 修复：限制x域到0.7，防止exp计算数值爆炸
				\addplot[domain=0:0.72, samples=100, ultra thick, blue] {0.01*(exp(x/0.026)-1)};
				% Reverse characteristic (small constant)
				\addplot[domain=-5:0, samples=10, ultra thick, blue] {-0.001};
				% Breakdown (draw as steep line near -5V)
				% 修复：修改坐标保证不超出ymin
				\draw[ultra thick, red, ->] (axis cs:-5,-0.9) -- (axis cs:-5,-0.1);
				\node[right] at (axis cs:0.6, 15) {\textbf{正向导通区}};
				\node[left] at (axis cs:-4, 2) {\textbf{反向截止区}};
				\node[right, red] at (axis cs:-5, -0.5) {\textbf{击穿区}};
				\draw[dashed, gray] (axis cs:0.7, 0) -- (axis cs:0.7, -1);
				\node[below] at (axis cs:0.7, -1) {$U_{ON}$};
			\end{axis}
		\end{tikzpicture}
	\end{center}
	在分析电路时，常根据精度需求选择模型：
	\begin{enumerate}
		\item \textbf{理想模型}：正偏视作短路（$U_D=0$），反偏视作开路。
		\item \textbf{恒压降模型}：硅管导通时 $U_D = 0.7\text{V}$，锗管 $U_D = 0.3\text{V}$。
		\item \textbf{小信号模型}：在直流工作点附近，将二极管等效为动态电阻 $r_d = \frac{U_T}{I_D}$。
	\end{enumerate}
	\subsubsection*{2.3.2 稳压管}
	稳压管工作在\textbf{反向击穿区}。
	\begin{tipbox}[形象理解稳压管：弹簧限位器]
		想象一面墙前装了一个弹簧限位器。当人撞向墙壁（电压升高）时，弹簧被压缩，吸收了大部分冲击力（电流急剧增大），而人受到的实际冲击力（两端电压）基本保持恒定。稳压管就是那个牺牲自己（增大电流）来保护电压恒定的限位器。
	\end{tipbox}
	\subsection*{2.4 双极型晶体管 (BJT)}
	\subsubsection*{2.4.1 结构与电流分配}
	BJT有三个区：发射区（重掺杂）、基区（极薄且掺杂低）、集电区。对应三个电极：发射极E、基极B、集电极C。分为NPN和PNP两种。
	\begin{notebox}[BJT放大的核心逻辑]
		在放大状态下（发射结正偏，集电结反偏），发射区发射大量电子，由于基区极薄且空穴极少，只有极少数电子在基区复合形成 $I_B$，绝大部分电子被集电结的强电场拉入集电区形成 $I_C$。这就是\textbf{小电流 $I_B$ 控制大电流 $I_C$} 的本质。
	\end{notebox}
	电流关系式：
	$$ I_C = \beta I_B, \quad I_E = I_B + I_C = (1+\beta)I_B \approx I_C $$
	\subsubsection*{2.4.2 BJT的特性曲线}
	以共射极接法为例：
	\begin{center}
		\begin{tikzpicture}
			% Input characteristics
			\begin{axis}[
				name=plot1,
				width=6cm, height=5cm,
				title={\textbf{输入特性曲线}},
				xlabel={$U_{BE}$ / V},
				ylabel={$I_B$ / $\mu$A},
				xmin=0, xmax=1, ymin=0, ymax=60,
				axis lines=left,
				thick,
				grid=major, grid style={dashed, gray!30},
				legend pos=north west,
				restrict y to domain=0:60 % 修复：防溢出
				]
				\addplot[domain=0:0.85, samples=50, ultra thick, blue] {0.01*(exp(x/0.026)-1)};
				\addlegendentry{$U_{CE} \ge 1V$}
			\end{axis}
			% Output characteristics
			\begin{axis}[
				name=plot2,
				at={($(plot1.east)+(3cm,0)$)}, anchor=west,
				width=7cm, height=5cm,
				title={\textbf{输出特性曲线}},
				xlabel={$U_{CE}$ / V},
				ylabel={$I_C$ / mA},
				xmin=0, xmax=12, ymin=0, ymax=6,
				axis lines=left,
				thick,
				grid=major, grid style={dashed, gray!30},
				]
				% I_B = 0
				\addplot[domain=0:12, samples=10, ultra thick, blue] {0.01};
				% I_B = 20uA
				\addplot[domain=0.3:12, samples=50, ultra thick, red] {1.0 + 0.01*x};
				% I_B = 40uA
				\addplot[domain=0.3:12, samples=50, ultra thick, orange] {2.0 + 0.02*x};
				% I_B = 60uA
				\addplot[domain=0.3:12, samples=50, ultra thick, green!60!black] {3.0 + 0.03*x};
				% Saturation region
				\fill[gray!20] (axis cs:0,0) rectangle (axis cs:0.7,6);
				\node at (axis cs:0.35, 5) {\textbf{饱和区}};
				% Active region
				\node at (axis cs:6, 5) {\textbf{放大区}};
				% Cutoff region
				\node at (axis cs:10, 0.5) {\textbf{截止区}};
			\end{axis}
		\end{tikzpicture}
	\end{center}
	三个工作区的判断：
	\begin{enumerate}
		\item \textbf{放大区}：发射结正偏，集电结反偏。$I_C = \beta I_B$，曲线水平。
		\item \textbf{饱和区}：发射结正偏，集电结正偏。$U_{CE} < U_{BE}$，$I_C < \beta I_B$，开关“闭合”。
		\item \textbf{截止区}：发射结反偏（或零偏），集电结反偏。$I_B=0, I_C \approx 0$，开关“断开”。
	\end{enumerate}
	\subsection*{2.5 场效应管 (FET)}
	FET是\textbf{电压控制器件}（$I_D$ 受 $U_{GS}$ 控制），输入阻抗极高。
	\begin{itemize}
		\item \textbf{结型 (JFET)}：利用PN结反向偏置改变耗尽层宽度来控制导电沟道。
		\item \textbf{绝缘栅 (MOSFET)}：分为增强型（无原始沟道，需开启电压 $U_{GS(th)}$）和耗尽型（有原始沟道）。
	\end{itemize}
	增强型NMOS的转移特性公式（在饱和区 $U_{DS} > U_{GS} - U_{GS(th)}$）：
	$$ I_D = I_{DO}\left(1 - \frac{U_{GS}}{U_{GS(th)}}\right)^2 $$
	其中 $I_{DO}$ 是 $U_{GS} = 2U_{GS(th)}$ 时的漏极电流。
	\begin{tipbox}[形象理解MOSFET：水闸门]
		栅极G就像水龙头上的旋钮，改变旋钮开度（$U_{GS}$）就能控制水流（$I_D$）大小。旋钮本身不接触水流（绝缘栅），因此不消耗水（输入电流极小）。
	\end{tipbox}
	\clearpage
	\section*{第3章 放大电路基础}
	\addcontentsline{toc}{section}{第3章 放大电路基础}
	\subsection*{3.1 为什么需要静态工作点 (Q点)？}
	如果直接把交流信号加在BJT上，信号的负半周会使发射结反偏，BJT截止，造成\textbf{严重失真}。
	\begin{warnbox}[核心原则：先直流，后交流]
		设置Q点就像在跷跷板中间加个支点，让交流信号在“导通区”内上下波动。直流是舞台，交流是演员，没有舞台，演员就无法表演。
	\end{warnbox}
	\subsection*{3.2 基本共射放大电路分析}
	以经典的阻容耦合共射放大电路为例。
	% 修复：重绘共射放大电路，确保电路拓扑逻辑正确
	\begin{center}
		\begin{circuitikz}[scale=1.0, transform shape, european]
			\draw (0,0) node[ground]{};
			\draw (0,6) node[above]{$V_{CC}$};
			% Rb
			\draw (0,6) to[R, l=$R_b$] (3,6) -- (3,4.5);
			% BJT
			\draw (3,4.5) node[npn, anchor=B](Q){};
			% Rc
			\draw (Q.C) to[R, l=$R_c$] (3,6);
			% Re
			\draw (Q.E) to[R, l=$R_e$] (3,0) node[ground]{};
			% Input coupling
			\draw (-2,4.5) to[C, l=$C_1$] (Q.B);
			\draw (-2,4.5) node[left]{$u_i$};
			% Output coupling
			\draw (Q.C) -- (4.5,3.5);
			\draw (4.5,3.5) to[C, l=$C_2$] (6.5,3.5);
			\draw (6.5,3.5) to[R, l=$R_L$] (6.5,0) node[ground]{};
			\draw (6.5,3.5) -- (7.5,3.5) node[right]{$u_o$};
		\end{circuitikz}
	\end{center}
	\subsubsection*{3.2.1 直流通路与静态分析}
	直流通路：电容开路，电感短路。
	\begin{derivebox}[估算法求Q点推导]
		1. 基极回路 (KVL)：
		$$ V_{CC} = I_B R_b + U_{BEQ} $$
		得出：$ I_B = \frac{V_{CC} - U_{BEQ}}{R_b} \approx \frac{V_{CC}}{R_b} $ (当 $V_{CC} \gg U_{BEQ}$时)
		2. 集电极回路：
		$$ I_C = \beta I_B $$
		3. 管压降 (KVL)：
		$$ U_{CE} = V_{CC} - I_C R_c $$
	\end{derivebox}
	\subsubsection*{3.2.2 图解法分析}
	图解法能直观展示Q点的位置与波形失真的关系。
	\begin{center}
		\begin{tikzpicture}
			\begin{axis}[
				width=10cm, height=8cm,
				xlabel={$U_{CE}$ / V},
				ylabel={$I_C$ / mA},
				xmin=0, xmax=15, ymin=0, ymax=5,
				axis lines=left,
				thick,
				grid=major, grid style={dashed, gray!30},
				xtick={0, 2.5, 5, 10, 12, 15},
				ytick={0, 1, 2, 3, 4, 5},
				]
				% DC Load line: I_C = (Vcc - U_CE) / Rc. Vcc=12, Rc=3k -> I_C = 4 - U_CE/3
				\addplot[domain=0:12, ultra thick, blue] {4 - x/3};
				\node[above, blue] at (axis cs:11, 0.5) {\textbf{直流负载线}};
				% Output characteristic curves
				\addplot[domain=0.3:15, samples=50, thick, red] {0.5 + 0.02*x};
				\addplot[domain=0.3:15, samples=50, thick, red] {1.5 + 0.02*x};
				\addplot[domain=0.3:15, samples=50, thick, red] {2.5 + 0.02*x};
				\addplot[domain=0.3:15, samples=50, thick, red] {3.5 + 0.02*x};
				% Q Point
				\addplot[only marks, mark=*, mark size=3pt, black] coordinates {(6, 2)};
				\node[above right, font=\bfseries] at (axis cs:6, 2) {Q点 ($6V, 2mA$)};
				% AC Load line (Assuming Rc//Rl = 1.5k, Ic swing +/- 1mA around Q)
				% Slope = -1/1.5 = -0.66. Passes through Q(6,2). Ic_max=3, Uce_max=7.5
				\addplot[domain=1.5:10.5, ultra thick, green!60!black, dashed] {2 - (2/3)*(x-6)};
				\node[below left, green!60!black] at (axis cs:10, 0.6) {\textbf{交流负载线}};
				% Saturation and Cutoff marks
				\draw[->, thick, orange] (axis cs:1, 3.5) -- (axis cs:1, 4.5) node[above] {饱和失真};
				\draw[->, thick, purple] (axis cs:12, 0.5) -- (axis cs:13, 0.5) node[right] {截止失真};
			\end{axis}
		\end{tikzpicture}
	\end{center}
	\begin{warnbox}[失真分析]
		\begin{itemize}
			\item \textbf{饱和失真}：Q点过高 $\to$ 信号正半周进入饱和区 $\to$ 输出波形底部被削平（共射电路中）。
			\item \textbf{截止失真}：Q点过低 $\to$ 信号负半周进入截止区 $\to$ 输出波形顶部被削平。
		\end{itemize}
	\end{warnbox}
	\subsubsection*{3.2.3 微变等效电路法 (小信号模型)}
	交流通路：大容量电容短路，直流电源置地（内阻为0）。
	将BJT在Q点处等效为线性二端口网络：
	\begin{itemize}
		\item 输入端等效为动态电阻 $r_{be}$：
		$$ r_{be} = r_{bb'} + (1+\beta)\frac{U_T}{I_{EQ}} \approx 300\Omega + (1+\beta)\frac{26\text{mV}}{I_{CQ}\text{(mA)}} $$
		\item 输出端等效为受控电流源 $\beta i_b$。
	\end{itemize}
	\begin{derivebox}[共射放大电路交流参数推导]
		1. 电压放大倍数 $A_u$：
		$$ A_u = \frac{u_o}{u_i} = \frac{-\beta i_b (R_c \parallel R_L)}{i_b r_{be}} = -\frac{\beta R'_L}{r_{be}} $$
		\textbf{注意}：负号表示输出与输入反相，这是共射电路的标志性特征！
		2. 输入电阻 $R_i$：
		$$ R_i = R_b \parallel r_{be} \approx r_{be} $$
		3. 输出电阻 $R_o$：
		令 $u_i=0$，保留负载，在输出端加电压 $u$ 求电流 $i$。受控源 $\beta i_b=0$ 相当于开路。
		$$ R_o \approx R_c $$
	\end{derivebox}
	\subsection*{3.3 三种基本组态比较}
	\begin{center}
		\renewcommand{\arraystretch}{1.8}
		\begin{tabular}{|c|c|c|c|}
			\hline
			\textbf{特性} & \textbf{共射 (CE)} & \textbf{共集 (CC/射极跟随器)} & \textbf{共基 (CB)} \\
			\hline
			$A_u$ & 大（反相） & 接近1（同相） & 大（同相） \\
			\hline
			$R_i$ & 中 & 大（隔离信号源） & 小 \\
			\hline
			$R_o$ & 中 & 小（驱动负载） & 大 \\
			\hline
			\textbf{形象比喻} & 核心主力放大 & 阻抗变换缓冲器 & 高频宽带放大器 \\
			\hline
		\end{tabular}
	\end{center}
	\clearpage
	\section*{第4章 反馈放大电路}
	\addcontentsline{toc}{section}{第4章 反馈放大电路}
	\subsection*{4.1 反馈的基本概念}
	\textbf{反馈}：将输出量（电压或电流）通过一定网络送回输入回路，影响输入量的过程。
	\begin{itemize}
		\item \textbf{正反馈}：反馈信号增强原输入信号，容易振荡（如话筒靠近音响啸叫）。
		\item \textbf{负反馈}：反馈信号削弱原输入信号，牺牲增益换取稳定性。
	\end{itemize}
	\begin{center}
		\begin{tikzpicture}[node distance=2.5cm, >=stealth, block/.style={draw, fill=blue!20, minimum width=2cm, minimum height=1cm, thick}]
			\node[block] (A) {$A$ (开环放大)};
			% 修复：使用 positioning 库的 below=of 语法
			\node[block, below=1.5cm of A] (F) {$F$ (反馈网络)};
			\node[circle, draw, thick, left=2cm of A] (sum) {};
			\draw[->, thick] (sum) -- node[above] {$X_{id}$} (A);
			\draw[->, thick] ([xshift=-1.5cm]sum.west) node[left] {$X_i$} -- (sum);
			\draw[->, thick] (A) -- node[above] {} ++(3,0) node[right] {$X_o$ (输出)};
			\draw[->, thick] (A.east) ++(1.5,0) |- (F);
			\draw[->, thick] (F) -| (sum) node[pos=0.95, left] {$-$} node[pos=0.95, right] {$X_f$};
			\node at (sum) {$\Sigma$};
		\end{tikzpicture}
	\end{center}
	基本方程：$\displaystyle A_f = \frac{X_o}{X_i} = \frac{A}{1+AF}$。其中 $AF$ 为环路增益，$|1+AF|$ 为反馈深度。
	\subsection*{4.2 负反馈的四种组态判断}
	判断组态是考研和期末考的重点，遵循“两步法”。
	\begin{notebox}[组态判断秘籍]
		\textbf{第一步：判断是电压还是电流反馈（看输出）}
		将输出端短路（$u_o=0$），若反馈信号消失，则为\textbf{电压反馈}；若依然存在，则为\textbf{电流反馈}。
		\textbf{形象记忆}：电压反馈是“测电压”，并联在输出端；电流反馈是“测电流”，串联在输出端。
		\textbf{第二步：判断是串联还是并联反馈（看输入）}
		若反馈信号与输入信号接在\textbf{同一个极}（如都在基极），则为\textbf{并联反馈}；接在\textbf{不同极}（如输入在基极，反馈在发射极），则为\textbf{串联反馈}。
		\textbf{形象记忆}：并联是电流相加减，串联是电压相加减。
	\end{notebox}
	四种组态及其对信号源的要求：
	\begin{enumerate}
		\item \textbf{电压串联}：稳定输出电压，增大输入电阻。要求信号源内阻小（恒压源）。
		\item \textbf{电压并联}：稳定输出电压，减小输入电阻。要求信号源内阻大（恒流源）。
		\item \textbf{电流串联}：稳定输出电流，增大输入电阻。要求信号源内阻小。
		\item \textbf{电流并联}：稳定输出电流，减小输入电阻。要求信号源内阻大。
	\end{enumerate}
	\subsection*{4.3 负反馈对放大电路性能的影响}
	\begin{tipbox}[负反馈的“四大功劳”]
		牺牲了增益 $A$ 降到 $A_f$，换取了以下性能的提升（均与反馈深度 $1+AF$ 相关）：
		\begin{enumerate}
			\item \textbf{稳}：提高增益的稳定性。$\frac{dA_f}{A_f} = \frac{1}{1+AF}\frac{dA}{A}$。
			\item \textbf{宽}：展宽频带。$BW_f = (1+AF)BW$。犹如把高耸的山峰压成平缓的丘陵。
			\item \textbf{小}：减小非线性失真和内部噪声。只能减小环路内部的失真，且输出幅度不变。
			\item \textbf{调}：对输入输出电阻的调节。
			\begin{itemize}
				\item 串联负反馈使输入电阻增大 $(1+AF)$ 倍；
				\item 并联负反馈使输入电阻减小 $(1+AF)$ 倍；
				\item 电压负反馈使输出电阻减小 $(1+AF)$ 倍（恒压源特性）；
				\item 电流负反馈使输出电阻增大 $(1+AF)$ 倍（恒流源特性）。
			\end{itemize}
		\end{enumerate}
	\end{tipbox}
	\subsection*{4.4 深度负反馈下的估算}
	当 $|1+AF| \gg 1$ 时，$A_f \approx \frac{1}{F}$。
	这意味着闭环增益仅由反馈网络（通常是无源电阻网络）决定，与有源器件参数无关！
	推导引出两条黄金法则：
	$$ X_i \approx X_f \Rightarrow \begin{cases} \text{串联反馈: } u_i \approx u_f & (\text{虚短}) \\ \text{并联反馈: } i_i \approx i_f & (\text{虚断}) \end{cases} $$
	\clearpage
	\section*{第5章 集成运算放大器}
	\addcontentsline{toc}{section}{第5章 集成运算放大器}
	集成运放是一个高增益、高输入电阻、低输出电阻的直接耦合多级放大电路。
	\subsection*{5.1 差分放大电路}
	直接耦合放大器的最大痛点是\textbf{零点漂移}（温漂）。差放是抑制零漂的利器。
	\begin{notebox}[形象理解差放：双胞胎兄弟]
		差放由两个完全对称的“单管放大器”（双胞胎）组成。温度变化相当于给两人同时加了相同的输入（共模信号）。因为两人本领一样，反应也完全一样，在双端输出时，两人的反应相互抵消，输出为0。
		而对有用的信号，以相反极性加给两人（差模信号），一升一降，输出就会产生巨大差异，从而实现“放大差模，抑制共模”。
	\end{notebox}
	% 修复：改善差分放大电路绘图，去除不规范的 node[vee]
	\begin{center}
		\begin{circuitikz}[scale=0.9, transform shape, european]
			\draw (0,0) node[ground]{};
			\draw (-2,4) node[npn](Q1){T1};
			\draw (2,4) node[npn, xscale=-1](Q2){T2};
			% VCC
			\draw (0,6) node[above]{$V_{CC}$};
			\draw (-2,6) to[R, l=$R_c$] (-2,4.5);
			\draw (2,6) to[R, l_=$R_c$] (2,4.5);
			\draw (-2,6) -- (2,6);
			% Tail Resistor
			\draw (-2,3.5) -- (-2,2);
			\draw (2,3.5) -- (2,2);
			\draw (-2,2) -- (2,2);
			\draw (0,2) to[R, l=$R_e$] (0,0);
			% Inputs
			\draw (Q1.B) to[short] ++(-1,0) node[left]{$u_{i1}$};
			\draw (Q2.B) to[short] ++(1,0) node[right]{$u_{i2}$};
			% Outputs
			\draw (-2,6) -- (-2,6.5);
			\draw (2,6) -- (2,6.5);
			\draw[->, thick] (-2,6.5) -- (-3,7) node[above]{$u_{oc1}$};
			\draw[->, thick] (2,6.5) -- (3,7) node[above]{$u_{oc2}$};
		\end{circuitikz}
	\end{center}
	核心指标——\textbf{共模抑制比}：
	$$ K_{CMR} = \left| \frac{A_{ud}}{A_{uc}} \right| $$
	理想对称双端输出时，$A_{uc} = 0$，$K_{CMR} \to \infty$。发射极电阻 $R_e$ 的作用是引入共模负反馈，进一步抑制单端输出的共模信号。
	\subsection*{5.2 理想运算放大器及其两条“黄金法则”}
	理想运放参数：开环增益 $A_{od} \to \infty$，输入电阻 $R_{id} \to \infty$，输出电阻 $R_o \to 0$，共模抑制比 $K_{CMR} \to \infty$。
	\begin{warnbox}[运放线性应用的两条黄金法则]
		运放工作在线性区（必须引入负反馈！）时：
		\begin{enumerate}
			\item \textbf{虚短}：$u_+ = u_-$。因为 $A_{od} \to \infty$ 而输出有限，两输入端电位逼近相等，似短路而非短路。
			\item \textbf{虚断}：$i_+ = i_- = 0$。因为输入电阻 $\to \infty$，两输入端几乎不取电流，似断路而非断路。
		\end{enumerate}
		这两条法则是分析一切运放线性应用电路的万能钥匙！
	\end{warnbox}
	\subsection*{5.3 基本运算电路}
	\subsubsection*{5.3.1 反相比例运算}
	信号从反相端输入，同相端接地。
	% 修复：使用规范美观的运放反馈电阻画法
	\begin{center}
		\begin{circuitikz}[european, scale=1.0]
			\draw (0,0) node[op amp] (opamp) {};
			\draw (opamp.+) to[short] ++(-0.5,0) node[ground]{};
			\draw (opamp.-) to[short] ++(-0.5,0) coordinate(inv);
			\draw (inv) to[R, l=$R_1$] ++(-2.5,0) node[left]{$u_i$};
			% 反馈电阻从反相端向上绕至输出端
			\draw (inv) to[R, l=$R_f$] ++(0, 2) coordinate(rftop) -| (opamp.out);
			\draw (opamp.out) to[short] ++(1,0) node[right]{$u_o$};
		\end{circuitikz}
	\end{center}
	\begin{derivebox}[推导过程]
		利用虚断：$i_1 = i_f$。
		利用虚短（虚地）：$u_- = u_+ = 0$。
		$$ i_1 = \frac{u_i - 0}{R_1}, \quad i_f = \frac{0 - u_o}{R_f} $$
		联立得：
		$$ \frac{u_i}{R_1} = -\frac{u_o}{R_f} \Rightarrow u_o = -\frac{R_f}{R_1} u_i $$
	\end{derivebox}
	\subsubsection*{5.3.2 同相比例运算}
	信号从同相端输入，反相端接反馈。
	\begin{derivebox}[推导过程]
		利用虚短：$u_- = u_+ = u_i$。
		利用虚断：$i_1 = i_f$。
		$$ i_1 = \frac{0 - u_i}{R_1}, \quad i_f = \frac{u_i - u_o}{R_f} $$
		联立得：
		$$ \frac{-u_i}{R_1} = \frac{u_i - u_o}{R_f} \Rightarrow u_o = \left(1 + \frac{R_f}{R_1}\right) u_i $$
		特例：当 $R_1 \to \infty$ (开路) 且 $R_f=0$ (短路) 时，$u_o = u_i$，称为\textbf{电压跟随器}，具有极高的输入电阻和极低的输出电阻，常作缓冲级。
	\end{derivebox}
	\subsubsection*{5.3.3 加减运算电路}
	\begin{itemize}
		\item \textbf{反相求和}：多路信号在反相端汇总，互不干扰。
		$$ u_o = -\left(\frac{R_f}{R_{11}}u_{i1} + \frac{R_f}{R_{12}}u_{i2}\right) $$
		\item \textbf{差分减法}：双端输入，利用叠加定理。
		$$ u_o = \left(1+\frac{R_f}{R_1}\right)\frac{R_3}{R_2+R_3}u_{i2} - \frac{R_f}{R_1}u_{i1} $$
		当电阻匹配满足 $\frac{R_f}{R_1} = \frac{R_3}{R_2}$ 时，$u_o = \frac{R_f}{R_1}(u_{i2} - u_{i1})$。
	\end{itemize}
	\subsubsection*{5.3.4 积分与微分运算}
	将电阻替换为电容，利用电容的充放电特性。
	\begin{itemize}
		\item \textbf{积分电路}：电容 $C$ 在反馈支路。
		$$ u_o = -\frac{1}{R_1 C}\int u_i dt $$
		\textbf{形象理解}：输入阶跃电压，电容缓慢充电，输出呈斜坡下降。积分器相当于电路的“记忆器”。
		\item \textbf{微分电路}：电容 $C$ 在输入支路。
		$$ u_o = -R_f C \frac{du_i}{dt} $$
		\textbf{形象理解}：只对变化的信号有反应，相当于电路的“速度计”。
	\end{itemize}
	% 修复：重绘积分微分波形图，修复坐标系超限问题
	\begin{center}
		\begin{tikzpicture}
			\begin{axis}[
				width=10cm, height=6cm,
				xlabel={时间 $t$},
				ylabel={电压 $u$},
				xmin=0, xmax=5, ymin=-2, ymax=3,
				axis lines=middle,
				thick,
				grid=major, grid style={dashed, gray!30},
				legend pos=south west,
				xticklabels={}, yticklabels={}
				]
				% Step input
				\addplot[domain=0:5, ultra thick, blue] {2};
				\addlegendentry{输入阶跃 $u_i$}
				% Integral output (ramp)
				\addplot[domain=0:5, ultra thick, red] {-0.5*x};
				\addlegendentry{积分输出 $u_o$ (斜坡)}
				% Differential output (impulse approx)
				\addplot[ultra thick, orange, const plot] coordinates {(0,0) (0.01,2.5) (0.05,0) (5,0)};
				\addlegendentry{微分输出 $u_o$ (冲击)}
			\end{axis}
		\end{tikzpicture}
	\end{center}
	\vspace{2em}
	\begin{center}
		\Large\textbf{第1$\sim$5章 模拟电路核心复习完毕！}
	\end{center}
\end{document}